\begin{proposition}[Renormalisation freedom of $Q_1$, $S_1$ and $K_1$]
\label{prop:c11-freedom}
Under the changes of Proposition \ref{prop:c8-laws} the components change as follows; $\{S^{(1)}(x),-\}$ is the classical transformation along $x$.
\begin{itemize}
\item (i) A field redefinition $F_{\mathfrak{psl}(4|4)}\mapsto F_{\mathfrak{psl}(4|4)}+\hbar s\phi$ changes the component along $x$ by $-\{(s\phi)^{(1)}(x),-\}$, a derivation of the letters: it changes the one-body kernels and leaves the two-body kernels unchanged. For the field countings $\phi^{(0)}=\sum_fy_f\int(f,f^*)$ the change of $Q_1$ multiplies each arrow $f\to f'$ of the classical transformation by $y_{f'}-y_f$; for $\phi^{(0)}=t\,S^{(1)}(\Delta)$ it is a multiple of the classical transformation proportional to $\lambda_x$.
\item (ii) The coupling direction leaves $Q_1$, $S_1$, $K_1$ unchanged. The outer direction $G_B$ changes $Q_1$ along $S_+\otimes W$ by $-a$ and along $S_-\otimes W^*$ by $+a$ times the classical transformation, $S_1$ along $S_+\otimes W^*$ by $+a$ and along $S_-\otimes W$ by $-a$ times it, and leaves $K_1$, $D_1$, $P_1$, $M_1$, $R_1$ unchanged.
\item (iii) The free directions of $\hbar F_1^{(1)}(x)$ change the one-body kernels: along $\Sigma$ by Proposition \ref{prop:c11-freexi} (iii); along $\Delta$ and $P_b$ through the sector coefficients of Definition \ref{def:c11-free}, which enter $D_1$, $P_1$ and, through $P_1$, the relations of Proposition \ref{prop:c11-S} and Corollary \ref{cor:c11-K}; along $K_b$ through the $s$-closed directions of degree $-1$ quadratic in the Killing data. $\hbar F_1^{(0)}$ enters the one-body kernels through \eqref{eq:c11-arity1xi} and through the trees with $S^{(1)}(x)$ (Remark \ref{rem:c10-scheme1} (a), (b)).
\item (iv) A finite part $\psi'$ of the operators changes the component along $x$ by $[\{S^{(1)}(x),-\},\psi']$, on $O$ by $\{S^{(1)}(x),\psi'(O)\}-\psi'(\{S^{(1)}(x),O\})$. For $\psi'=\tfrac a2\{S^{(1)}(\Delta),-\}$ this is $\tfrac a2\lambda_x\{S^{(1)}(x),-\}$: $+a$, $-\tfrac a2$, $+\tfrac a2$ times the classical transformation along $K_b$, $\Sigma$, $\Sigma^*$ (Proposition \ref{prop:c9-scheme}), a change of the one-body kernels. A finite part $\psi'$ that vanishes on $O$ changes the component of $O$ along $x$ by $-\psi'(\{S^{(1)}(x),O\})$: the finite parts of the operators of the Weyl weight of $\{S^{(1)}(x),O\}$ enter the two-body kernels, as the finite part of Lemma \ref{lem:c11-cocycle} on the trace bilinear and the shift of the Konishi entry in Remark \ref{rem:c11-twist}.
\item (v) The cutoff shape enters the finite parts of the defect cycles (Proposition \ref{prop:c10-shape}, Lemma \ref{lem:c10-special}), of the special-edge cycles of $L_1^{(0)}$ and of the finite special vertices; since the counterterms are unique up to finite local terms [C07, Thm B], a change of the shape is a change of $\hbar F_1$ and of $\psi$, a combination of (i)--(iv).
\item (vi) At the values $\theta=\theta_\lambda=\theta_V=-\tfrac13$ of Theorem \ref{thm:c9-theta2}, in the translation-covariant scheme, $P_1\cong0$ on the families of Definition \ref{def:c11-families} (Corollary \ref{cor:c11-D1fam}), and the condition $P_1=0$ of Definition \ref{def:c8-radial} holds there with $\psi'=0$: the components along $K_b$, $\Sigma$ and $\Sigma^*$ are those of the scheme. For equal coefficients $\theta_P\ne\theta$ along $P_b$, the finite part with $a=(\theta-\theta_P)$ units gives $P_1\cong0$ and changes the components along $K_b$, $\Sigma$, $\Sigma^*$ by $a$, $-\tfrac a2$, $\tfrac a2$ times the classical transformation (Corollary \ref{cor:c11-theta} (iv)).
\item (vii) Under (i), (iv) and the coupling direction the representation of $\mathfrak{psl}(4|4)$ on the cohomology of $\Lambda^{(0)}$ changes by an isomorphism, under $G_B$ by the automorphism $e^{-a\hbar B}$, and the spectrum of $D_1$ in a fixed classical Weyl weight is invariant (Theorem \ref{thm:c8-invariants}); the class $[Q^\xi_1]$ on the cohomology of $Q^\xi_0$, its kernel, image and rank are invariant under (i), (ii) and (iv) (Theorem \ref{thm:c8-twist}).
\end{itemize}
\end{proposition}
